\documentclass[a5paper,twoside]{ltjsarticle}
\usepackage{fancyhdr}
\pagestyle{fancy}
\fancyhf{}
\fancyhead[LE,RO]{柱の文字}
\fancyfoot[C]{\thepage}
\begin{document}
\section{柱とノンブル}
朝から細かい雨が降っていた。傘を持たずに家を出たのは失敗だったと、駅までの坂道を歩きながら何度も思った。商店街の軒先を伝って歩けば少しは濡れずに済むが、それでも肩のあたりはすぐに冷たくなった。

電車は思ったより空いていた。窓際の席に座ると、外の景色が雨粒で滲んで見えた。友人は鞄から文庫本を取り出し、しおりを挟んだ頁を開いた。

例の場所というのが何を指すのか、私にはすぐにわかった。学生の頃によく通った川沿いの喫茶店である。店主は無口な人で、注文を聞く以外にはほとんど口をきかなかった。

朝から細かい雨が降っていた。傘を持たずに家を出たのは失敗だったと、駅までの坂道を歩きながら何度も思った。商店街の軒先を伝って歩けば少しは濡れずに済むが、それでも肩のあたりはすぐに冷たくなった。

電車は思ったより空いていた。窓際の席に座ると、外の景色が雨粒で滲んで見えた。友人は鞄から文庫本を取り出し、しおりを挟んだ頁を開いた。

例の場所というのが何を指すのか、私にはすぐにわかった。学生の頃によく通った川沿いの喫茶店である。店主は無口な人で、注文を聞く以外にはほとんど口をきかなかった。

朝から細かい雨が降っていた。傘を持たずに家を出たのは失敗だったと、駅までの坂道を歩きながら何度も思った。商店街の軒先を伝って歩けば少しは濡れずに済むが、それでも肩のあたりはすぐに冷たくなった。

電車は思ったより空いていた。窓際の席に座ると、外の景色が雨粒で滲んで見えた。友人は鞄から文庫本を取り出し、しおりを挟んだ頁を開いた。

\clearpage
\thispagestyle{empty}
この頁は柱もノンブルもない。例の場所というのが何を指すのか、私にはすぐにわかった。学生の頃によく通った川沿いの喫茶店である。店主は無口な人で、注文を聞く以外にはほとんど口をきかなかった。

\clearpage
最後の頁。朝から細かい雨が降っていた。傘を持たずに家を出たのは失敗だったと、駅までの坂道を歩きながら何度も思った。商店街の軒先を伝って歩けば少しは濡れずに済むが、それでも肩のあたりはすぐに冷たくなった。
\end{document}
